\section{Motivation}
\label{sec:motivation}

Blockchain technology has emerged as a transformative innovation, providing a secure, distributed ledger that enables trust among untrusted entities without relying on centralized authorities or intermediaries. Its defining properties, including immutability, transparency, auditability, autonomy, and resilience, position blockchain as one of the most disruptive digital technologies of the past decade \cite{sanka2021systematic}.

Initially conceived as the foundation of Bitcoin \cite{nakamoto2008bitcoin}, blockchain’s applications have rapidly expanded beyond cryptocurrencies. Ethereum \cite{buterin2014ethereum} introduced programmability through smart contracts, enabling decentralized applications (dApps) that support diverse use cases. Today, blockchain is being deployed in finance, insurance, supply chain management, healthcare, identity management, registry services, stock trading, the Internet of Things (IoT), and energy systems.

Blockchain adoption has grown exponentially across various sectors. However, as the user base and transaction volumes expand, the underlying networks face significant scalability challenges. Ethereum Layer-1, a leading smart contract platform, provides strong security and decentralization but remains constrained by limited throughput, network congestion, and severe fee volatility. Under typical conditions, Ethereum's Layer-1 processes approximately 15–25 transactions per second (TPS), far below the thousands of TPS required for global-scale adoption. Its sequential execution model and consensus mechanism traditionally limit the network capacity.

This fundamental limitation reflects the widely recognized blockchain trilemma. The trilemma states that decentralized systems must balance scalability, security, and decentralization, and that all three properties cannot be maximized simultaneously. Consequently, a robust scaling strategy is required to achieve high transaction throughput and low latency without compromising the security guarantees of the base layer.

Ethereum's current scaling direction relies heavily on Layer-2 (L2) rollups. Rollups move transaction execution off-chain, grouping many off-chain transactions into large batches, while continuing to use Layer-1 (L1) for settlement, verification, and data availability. This approach amortizes the fixed L1 submission costs across multiple transactions, effectively increasing throughput while retaining Ethereum as the ultimate security layer.

Within this rollup-centric design space, Zero-Knowledge (ZK) Rollups are particularly attractive. They bundle hundreds or thousands of off-chain transfers into a single transaction and submit cryptographic validity proofs (such as zk-SNARKs or zk-STARKs) that allow Layer-1 smart contracts to mathematically verify off-chain state transitions. This avoids the need for L1 to re-execute every transaction, preserving the base layer's security model while significantly reducing the computational burden on L1. Recently, the Dencun upgrade implemented EIP-4844 (Proto-Danksharding) to introduce blob-carrying transactions, intended to reduce L2 data availability costs. However, under EIP-4844, blobs are sold in fixed 128~KB chunks. For low-volume rollups, this creates a ``Blob Dilemma'': wait to fill the blob (increasing delay) or post it partially empty (increasing cost per transaction). This highlights the need to understand configuration trade-offs in real-time.
